\subsection{贪心与基础数学}
\ojpractice{CSP201703A}{分蛋糕：第10次，2017-03，第1题；正重量累积到阈值后立即分出，第一篇练习给解答}\ojpractice{CSP202312B}{因子化简：第32次，2023-12，第2题；反复试除并记录完整指数；第三节给解答}\ojpractice{CSP202305B}{矩阵运算：第30次，2023-05，第2题；用结合律先算小矩阵，降低乘法次数}
\subsubsection{贪心需要证明}
贪心是在当前状态下确定一个选择并且不回退。它成立需要证明：至少存在一个最优解包含这次选择。常用方法是交换论证，即把某个最优方案换成包含该选择的方案，而目标值不变差。找到几个正确样例不是证明。

例如按价值最大的物品优先装背包并不总是正确：容量4，物品的（体积，价值）分别为 $(3,5),(2,4),(2,4)$，先拿价值5的物品只能得到5，拿后两件可得8。基础01背包通常需要DP，不能仅因“选大的”直观就使用贪心。

\subsubsection{例题 合并石堆的最小代价}
有 n 堆石子，每次可以任意选两堆合并，代价等于两堆石子数之和。求合成一堆的最小总代价。本题\textbf{没有相邻限制}。$1\le n\le2\times10^5$，每堆石子数在1到 $10^9$ 之间。输入第一行 n，第二行各堆大小。
\begin{verbatim}
输入
4
1 2 3 4
输出
19
\end{verbatim}
先合并1和2，代价3；再合并3和3，代价6；最后合并4和6，代价10。总计19。

\textbf{为什么取最小两堆。}合并过程可以看成二叉树，每个原始堆是叶子；总代价等于每个叶子权值乘以其深度之和。在固定树形中，较小权值放到较深位置不会使代价增大。选择一对最深的兄弟叶子，可通过交换把最小的两个权值放到那里。因此存在一个最优方案先合并最小两堆；把它们缩成一堆后，剩下同类子问题。重复此选择即可。

用小根堆维护所有堆，每次取两个最小值再插入它们的和。时间 $O(n\log n)$，空间 $O(n)$。只有一堆时无需合并，答案为0。
\Code{09_merge.cpp}

\subsubsection{中位数最小化绝对距离和}
若在数轴上选位置 x，使 $\sum_i|a_i-x|$ 最小，排序后取中位数即可。把最左和最右配对：对任何位于两者之间的 x，两项距离和都是两端距离；离开该区间只会增大。对称配对后，所有区间的公共部分包含中位数。偶数个元素时，中间两数之间的任意位置都可达到最小值。

例如 $[1,2,9]$ 取 x=2，距离和为8；取平均数4，距离和为10。平均数与中位数解决的目标不同。

\subsubsection{质因数分解与溢出检查}
对正整数 x，从2开始试除。若 d 整除 x，反复除去并统计指数。循环条件可写 \texttt{d <= x/d}，避免计算 $d^2$ 溢出。循环结束后若 x>1，它是最后一个质因子。1没有质因子。
\begin{lstlisting}
// prime[] and exponent[] each need at least 64 elements.
int factorize(long long x, long long prime[], int exponent[]) {
    int total = 0;
    for (long long d = 2; d <= x / d; d++) {
        if (x % d != 0) continue;
        int count = 0;
        do { x /= d; count++; } while (x % d == 0);
        prime[total] = d; exponent[total++] = count;
    }
    if (x > 1) { prime[total] = x; exponent[total++] = 1; }
    return total;
}
\end{lstlisting}
最坏试除次数为 $O(\sqrt{x})$。如果要分解很多较大的数，应重新估计总复杂度，而不是只看一次分解。

\subsubsection{矩阵运算先看维度}
若 A 是 $n\times k$，B 是 $k\times m$，则 $C=AB$ 是 $n\times m$，有 $C_{ij}=\sum_{t=1}^k A_{it}B_{tj}$，朴素复杂度 $O(nkm)$。矩阵乘法一般不可交换。若题意出现转置或连乘，应先写出各中间矩阵的维度与计算量，再决定结合顺序。计算量可以因结合顺序变化而显著不同，但结果在精确算术下符合结合律。

\subsubsection{自测与解答}
\textbf{练习1。}三堆 $[2,2,8]$ 的最优代价是多少？\textbf{解答。}先合并2和2得到4，代价4；再与8合并代价12，总计16。若先合并2和8则为22。

\textbf{练习2。}分解360。\textbf{解答。}$360=2^3\cdot3^2\cdot5$。反复整除才能记录完整指数。

\textbf{练习3。}A为 $10\times100$，B为 $100\times5$，C为 $5\times50$，先算哪一对更省？\textbf{解答。}$(AB)C$ 的乘加规模为 $10\cdot100\cdot5+10\cdot5\cdot50=7500$；$A(BC)$ 为 $100\cdot5\cdot50+10\cdot100\cdot50=75000$，前者更省。
\newpage

\subsubsection{当前团队题单训练}
\training{1112208}{贪心}
先做合并果子、线段覆盖、纪念品分组、删数，再做道路、队伍接续、编码树和国王游戏。各题的题意、解答与关键代码见本篇练习题中的团队专项。
\practice{P1090}{合并果子：小根堆}
\practice{P1803}{线段覆盖：结束时间优先}
\practice{P1094}{纪念品分组：排序与双指针}
\practice{P1106}{删数问题：单调栈}
